% Macro for code snippets
\newcommand{\sconscode}[3]
{
  \begin{block}{}
  \lstinputlisting[firstline={#1},lastline={#2}]{#3}
  \end{block}
}

%
% Document
%
\title[]{Plotting in Madagascar}

\author[Karimi]{Parvaneh Karimi}
\institute[RSF-dev team]{Madagascar development team}
\date[Madagascar school in Austin] % (optional)
{July 20, 2012}

\setbeamercolor{alerted text}{fg=white}
\begin{frame}
  \titlepage
  % Hide progress bar in the footline
  \appendix
\end{frame}

\begin{frame}
  \frametitle{Outline}
   \hspace*{2cm}
   \begin{minipage}[t][3cm]{10cm}
    \tableofcontents
   \end{minipage}
\end{frame}

\section{Plotting infrastructure}

\begin{frame}
  \begin{block}{Principles}
    \begin{itemize}
      \item All plotting programs take RSF files as input through stdin and write binary Vplot to stdout
      \item Device-dependent rendering is handled by Vplot pens
      \item Madagascar SCons build workflows will handle all conversions automatically
    \end{itemize}
  \end{block}
  \begin{figure}
  \includegraphics[scale=0.45]{Fig/RSF.pdf}
  \end{figure}
\end{frame}

\begin{frame}
  \begin{block}{Command line}
    \$ \textless input.rsf plotting\_program [parameters] \textgreater output.vpl \\
    \$ \textless input.vpl vplot\_pen [parameters] \\
    \$ \textless input.rsf some\textbar processing\textbar workflow\textbar plotting\textbar pen\textgreater image
  \end{block}

  \begin{block}{Manual conversion}
    \$ vpconvert file.vpl file.jpg \\
    \$ vpconvert format=tiff Fig/*.vpl
  \end{block}


  \begin{block}{MADAGASCAR plotting programs}
\vfill
    \begin{minipage}{0.4\textwidth}
    \begin{itemize}
      \item sfgraph
      \item sfgrey
      \item sfcontour
      \item sfdots
      \item sfwiggle
      \item sfbox
    \end{itemize}
    \end{minipage}
    \begin{minipage}{0.4\textwidth}
    \begin{itemize}
       \item sfthplot
       \item sfgraph3
      \item sfgrey3
      \item sfgrey4
      \item sfcontour3
    \end{itemize}
    \end{minipage}
  \end{block}
\end{frame}

\begin{frame}
  \begin{block}{Pens}
  \begin{table}
    \begin{center}
     \begin{tabular}{|c|c|c|}
      \hline    Pen name   &  Dependency   &    Output          \\
      \hline & & \\[-1em]
      \hline     xtpen     &     X11       &     Screen         \\
      \hline     oglpen    &   OpenGL      &     Screen         \\
      \hline     jpegpen   &   libjpeg     &     JPEG           \\
      \hline     tiffpen   &   libtiff     &     TIFF           \\
      \hline     ppmpen    &   NetPBM      &     PPM            \\
      \hline     gdpen     & LibGD, FFmpeg &  PNG,JPEG,GIF,MPEG \\
      \hline     pngpen    &   Cairo       &      PNG           \\
      \hline     svgpen    &   Cairo       &      SVG           \\
      \hline     pdfpen    &   Cairo       &      PDF           \\
      \hline     pspen     &    None       &   Postscript (PS,EPS)      \\
      \hline     vppen     &    None       &   Vplot filter     \\
      \hline
    \end{tabular}
   \end{center}
  \end{table}
  \end{block}
\end{frame}

\begin{frame}
  \begin{figure}
  \includegraphics[scale=0.3]{Fig/xtglpens.pdf}
  \end{figure}
  \begin{block}{}
   RSFSRC \$ ./configure SFPEN=oglpen
  \end{block}
\end{frame}

\begin{frame}
  \begin{block}{SConstruct functions}
   \lstinputlisting{plotmeta.txt}
  \end{block}
  \begin{block}{Command line}
   \$ scons \textless target\textgreater.view \# View result on the screen \\
   \$ scons \textless target\textgreater.print \# Send to a printer \\
   \$ scons \textless target\textgreater.lock \# Install figure into a separate location \\
   \$ scons \textless target\textgreater.test \# Compare current figure with the ``locked'' one
  \end{block}
  \begin{block}{Environment variables}
   {\bf RSFFIGS} - location for installation of figures \\
   {\bf PSPRINTER} - destination for lp command
  \end{block}
\end{frame}

\section{Examples}

\begin{frame}
  \sconscode{4}{5}{plot/SConstruct}
  \begin{figure}
  \includegraphics[scale=0.25]{plot/Fig/sinc.pdf}
  \end{figure}
\end{frame}

\begin{frame}
  \sconscode{8}{11}{plot/SConstruct}
  \begin{figure}
  \includegraphics[scale=0.25]{plot/Fig/sqline.pdf}
  \end{figure}
\end{frame}

\begin{frame}
  \sconscode{12}{15}{plot/SConstruct}
  \begin{figure}
  \includegraphics[scale=0.25]{plot/Fig/msincs.pdf}
  \end{figure}
\end{frame}

\begin{frame}
  \sconscode{16}{17}{plot/SConstruct}
  \begin{figure}
  \includegraphics[scale=0.25]{plot/Fig/mpsincs.pdf}
  \end{figure}
\end{frame}

\begin{frame}
  \sconscode{86}{91}{plot/SConstruct}
  \begin{figure}
  \includegraphics[scale=0.25]{plot/Fig/dots.pdf}
  \end{figure}
\end{frame}

\begin{frame}
  \sconscode{24}{27}{plot/SConstruct}
  \begin{figure}
  \includegraphics[scale=0.25]{plot/Fig/mbmincs.pdf}
  \end{figure}
\end{frame}

\begin{frame}
  \sconscode{30}{33}{plot/SConstruct}
  \begin{figure}
  \includegraphics[scale=0.25]{plot/Fig/sqsurf.pdf}
  \end{figure}
\end{frame}

\begin{frame}
  \sconscode{34}{34}{plot/SConstruct}
  \begin{figure}
  \includegraphics[scale=0.25]{plot/Fig/csqsurf.pdf}
  \end{figure}
\end{frame}

\begin{frame}
  \sconscode{37}{40}{plot/SConstruct}
  \begin{figure}
  \includegraphics[scale=0.25]{plot/Fig/csurf.pdf}
  \end{figure}
\end{frame}

\begin{frame}
  \sconscode{48}{51}{plot/SConstruct}
  \begin{figure}
  \includegraphics[scale=0.25]{plot/Fig/sgath.pdf}
  \end{figure}
\end{frame}

\begin{frame}
  \sconscode{54}{58}{plot/SConstruct}
  \begin{figure}
  \includegraphics[scale=0.25]{plot/Fig/fganom.pdf}
  \end{figure}
\end{frame}

\begin{frame}
  \sconscode{61}{63}{plot/SConstruct}
  \begin{figure}
  \includegraphics[scale=0.25]{plot/Fig/ganom3.pdf}
  \end{figure}
\end{frame}

\begin{frame}
  \sconscode{68}{73}{plot/SConstruct}
  \begin{figure}
  \includegraphics[scale=0.25]{plot/Fig/sganom.pdf}
  \end{figure}
\end{frame}

%\begin{frame}
%  \sconscode{74}{75}{plot/SConstruct}
%  \begin{figure}
%  \includegraphics[scale=0.25]{plot/Fig/sqplsin.pdf}
%  \end{figure}
%\end{frame}

\begin{frame}
  \begin{block}{Combinations}
    Movie \\
    Overlay \\

     \begin{figure}
     \includegraphics[scale=0.25]{plot/Fig/ganom-7.pdf}
     \end{figure}


  \end{block}
\end{frame}

\begin{frame}
  \begin{block}{Combinations}
    SideBySideAniso \\

     \begin{figure}
     \includegraphics[scale=0.25]{plot/Fig/sganom.pdf}
     \end{figure}

  \end{block}
\end{frame}

\begin{frame}
  \begin{block}{Combinations}
    SideBySideIso \\

     \begin{figure}
     \includegraphics[scale=0.25]{plot/Fig/ganom2.pdf}
     \end{figure}

  \end{block}
\end{frame}

%\begin{frame}
% \begin{block}{Combinations}
%   OverUnderAniso \\

%     \begin{figure}
%     \includegraphics[scale=0.25]{plot/Fig/ganom-3.pdf}
%     \end{figure}

%  \end{block}
%\end{frame}
    

%\begin{frame}
%  \begin{block}{Combinations}
%    OverUnderIso \\

 %    \begin{figure}
 %    \includegraphics[scale=0.25]{plot/Fig/ganom-5.pdf}
 %    \end{figure}


%  \end{block}
%\end{frame}
    
%\begin{frame}
%  \begin{block}{Combinations}
%    TwoRows \\
%
%     \begin{figure}
%     \includegraphics[scale=0.25]{plot/Fig/ganom-5.pdf}
 %    \end{figure}

%  \end{block}
%\end{frame}
        
%\begin{frame}
%  \begin{block}{Combinations}
%    TwoColumns \\

%     \begin{figure}
 %    \includegraphics[scale=0.25]{plot/Fig/ganom-6.pdf}
%     \end{figure}

%  \end{block}
%\end{frame}   


\begin{frame}   
  \begin{block}{Command line help for common parameters}
    \$ sfdoc stdplot
  \end{block}
  \begin{block}{Use functions to reduce code redundancy}
%   \vspace{-0.5cm}
   \lstinputlisting{plotdef.txt}
  \end{block}
\end{frame}

\section{Plots in reproducible documents}

\begin{frame}
%  \vspace{0.5cm}
  \begin{figure}
  \includegraphics[scale=0.55]{Fig/pstruct.pdf}
  \caption{Directory structure for reproducible documents}
  \end{figure}
\end{frame}

\begin{frame}
  \begin{block}{paper/sample1/SConstruct}
  \lstinputlisting{sample1/SConstruct}
  \end{block}
  \begin{block}{paper/sample2/SConstruct}
  \lstinputlisting{sample2/SConstruct}
  \end{block}
\end{frame}

\begin{frame}
  \begin{block}{paper/sample.tex}
  \lstinputlisting[language=TeX,
                   morekeywords={begin,end,plot,title,author,
                                 subsection,ref,inputdir}]
                   {sample.tex}
  \end{block}
\end{frame}

\begin{frame}
  \begin{block}{paper/SConstruct}
  \lstinputlisting[language=python,firstline=1,lastline=5]{SConstruct}
%  \vspace{3mm}
%  \lstinputlisting[firstline=5,lastline=5]{SConstruct}
  \end{block}
  \begin{block}{Command line}
    \$ sftour scons lock \\
    \$ scons sample.pdf \\
    \$ scons sample.html \\
    \$ scons sample.install
  \end{block}
\end{frame}

\begin{frame}
  \begin{columns}
    \begin{column}{5.5cm}
      \begin{block}{}
      \includegraphics[page=1,scale=0.25]{sample.pdf}
      \end{block}
    \end{column}
    \begin{column}{5.5cm}
      \begin{block}{}
      \includegraphics[page=2,scale=0.25]{sample.pdf}
      \end{block}
    \end{column}
  \end{columns}
\end{frame}

%\begin{frame}
 % \begin{block}{Takeaway message}
  %  You do not have to convert anything manually, if you follow the workflow above.
  %\end{block}
%\end{frame}

